\documentclass[10pt]{article}

\usepackage[T1]{fontenc}
\usepackage{lmodern}
\usepackage{amsmath,amssymb,amsthm}
\usepackage[margin=1in]{geometry}
\usepackage{microtype}
\usepackage{xurl}
\usepackage[hidelinks]{hyperref}

\newtheorem{theorem}{Theorem}[section]
\newtheorem{lemma}[theorem]{Lemma}
\newtheorem{proposition}[theorem]{Proposition}
\newtheorem{corollary}[theorem]{Corollary}
\theoremstyle{definition}
\newtheorem{definition}[theorem]{Definition}
\allowdisplaybreaks

\title{A near-quadratic lower bound for sets with no unique sums}
\author{Working draft for proof audit}
\date{July 21, 2026}

\begin{document}

\maketitle

\begin{abstract}
For a prime $p$, let $m(p)$ be the least cardinality of a set
$A\subseteq\mathbb F_p$, $|A|\ge2$, for which every element of $A+A$ has at
least two distinct representations by unordered pairs from $A$.  We prove
\[
    m(p)\gg
    \left(\frac{\log p}{\log\log p}\right)^2.
\]
Together with the known construction $m(p)\ll(\log p)^2$, this gives
$m(p)=(\log p)^{2+o(1)}$.  The argument associates an integral collision
lattice to an inclusion-minimal example.  A clean-rectangle determinant
bound supplies a short relation between every pair of light blocks, and a
batched contraction reduces the lattice to
$O(\sqrt n+n/\log p)$ core coordinates while keeping all coefficients
polynomial in $n$.
\end{abstract}

\section{Statement and collision lattice}

\begin{definition}
Let $A$ be a finite subset of an abelian group.  An \emph{unordered
representation} of $s\in A+A$ is a two-element multiset $\{a,b\}$ with
$a,b\in A$, where $a=b$ is allowed, such that
\[
    s=a+b.
\]
Thus $\{a,b\}$ and $\{b,a\}$ are the same representation.  We say that
$A$ \emph{has no unique sum} if every $s\in A+A$ has at least two distinct
unordered representations.

Equivalently, for every two-element multiset $\{a,b\}$ from $A$, there is a
different multiset $\{c,d\}\ne\{a,b\}$ from $A$ such that
\[
    a+b=c+d.
\]
\end{definition}

Thus Problem 27 in Green's list asks for the asymptotic order of
\[
    m(p):=\min\bigl\{|A|:A\subseteq\mathbb F_p,\ |A|\ge2,\
                         A\text{ has no unique sum}\bigr\}.
\]
Bedert proved the previously strongest lower bound and an
$O((\log p)^2)$ upper bound~\cite{bedert}.  His construction, together with
the observation that $B+B$ has no unique sum whenever $B$ is balanced,
gives
\[
    m(p)\le\left(\frac12+o(1)\right)(\log_2p)^2.
\]

\begin{theorem}
There is an absolute constant $c>0$ such that, for every sufficiently large
prime $p$,
\[
    m(p)\ge
    c\left(\frac{\log p}{\log\log p}\right)^2.
\]
\end{theorem}

We begin the proof by introducing the collision lattice associated with an
inclusion-minimal set having no unique sum.  Fix a set
\[
    A=\{a_1,\ldots,a_n\}\subseteq\mathbb F_p
\]
that is inclusion-minimal among sets having no unique sum.  Every set having
no unique sum contains such an inclusion-minimal subset, so it suffices to
bound $n$.

Put
\[
    \Omega:=\{(i,j):1\le i\le j\le n\}.
\]
We regard each $(i,j)\in\Omega$ as an unordered index pair; the condition
$i\le j$ chooses a unique ordering, and pairs with $i=j$ are allowed.  For
$t=(i,j)\in\Omega$, put
\[
    u_t:=e_i+e_j\in\mathbb Z^n.
\]
In particular, $u_{(i,i)}=2e_i$.

Define the sum map
\[
    \sigma:\Omega\longrightarrow\mathbb F_p,
    \qquad
    \sigma(i,j)=a_i+a_j.
\]
Partition $\Omega$ into the fibers of $\sigma$; thus, two index pairs lie in
the same fiber exactly when they represent the same sum.  Since $A$ has no
unique sum, every fiber contains at least two pairs.  On every fiber choose
one cyclic permutation, and write $\mu:\Omega\to\Omega$ for the union of
these cycles.  In particular, $\mu$ has no fixed point.  Define
\[
    \rho_t:=u_t-u_{\mu(t)},
    \qquad
    \Lambda:=\operatorname{span}_{\mathbb Z}
                  \{\rho_t:t\in\Omega\}\subseteq\mathbb Z^n.
\]

\begin{lemma}[Rank and modular defect]
The lattice $\Lambda$ has rational rank $n-1$.  Modulo $p$, its row rank is
at most $n-2$.  Consequently every $(n-1)$-row matrix with rows in
$\Lambda$ has all its maximal minors divisible by $p$.
\end{lemma}

\begin{proof}
Let $x=(x_1,\ldots,x_n)\in\mathbb Q^n$ be orthogonal to every $\rho_t$.
Then $x_i+x_j$ is constant as $t=(i,j)$ runs around any one
$\mu$-cycle.  Suppose that $x$ is not constant, and let $M$ be the nonempty
proper set of indices on which $x_i$ is maximal.  If $t=(i,j)$ has both
$i,j\in M$, then $x_i+x_j$ is the largest possible pair value.  Equality
\[
    u_t\cdot x=u_{\mu(t)}\cdot x
\]
therefore forces both endpoints of $\mu(t)$ to lie in $M$.  Hence $\mu$
maps the finite set of pairs $(i,j)\in\Omega$ with $i,j\in M$ into itself
and therefore onto itself.  Since $\mu(t)\ne t$, every such index pair has a
distinct equal-sum mate whose two endpoints also lie in $M$.  The set $M$
is not a singleton, because its loop would otherwise have no distinct
internal mate.  Thus $\{a_i:i\in M\}$ is a proper subset of $A$ having no
unique sum, contrary to inclusion-minimality.  It follows that the rational
kernel consists exactly of the constant vectors, and hence
\[
    \operatorname{rank}_{\mathbb Q}\Lambda=n-1.
\]

Over $\mathbb F_p$, both
\[
    \mathbf 1=(1,\ldots,1)
    \qquad\text{and}\qquad
    a=(a_1,\ldots,a_n)
\]
are orthogonal to every $\rho_t$.  They are linearly independent because
the elements of $A$ are distinct and $n\ge2$.  The modular row rank is
therefore at most $n-2$.  Any maximal minor of any $(n-1)$-row matrix whose
rows lie in $\Lambda$ vanishes modulo $p$, which proves the final assertion.
\end{proof}

\section{Clean rectangles}

For disjoint nonempty $X,Y\subseteq[n]$, write $X\times Y$ for the set of
unordered index pairs having one index in $X$ and the other in $Y$.  We call
this the \emph{rectangle determined by $X$ and $Y$}.

\begin{lemma}[Clean rectangle]
If
\[
    \mu(X\times Y)\subseteq X\times Y,
\]
then
\[
    |X|+|Y|\ge\log_2p+2.
\]
\end{lemma}

\begin{proof}
Since $\mu$ is a bijection and $X\times Y$ is finite, the displayed
inclusion is equality.  Choose nonempty $X'\subseteq X$ and
$Y'\subseteq Y$ so that $X'\times Y'$ is invariant and
$|X'|+|Y'|$ is minimal.  Put
\[
    x=|X'|,
    \qquad
    y=|Y'|,
\]
and consider the rows $\rho_t$, $t\in X'\times Y'$, restricted to the
$x+y$ columns $X'\cup Y'$.

We first claim that their rational rank is $x+y-2$.  Their kernel contains
the two vectors that are constant on one side and zero on the other.
Conversely, let $(\alpha_i)_{i\in X'}$ and $(\beta_j)_{j\in Y'}$ lie in the
kernel.  Thus, whenever the index pair with endpoints $i,j$ is mapped by
$\mu$ to the index pair with endpoints $i',j'$,
\[
    \alpha_i+\beta_j=\alpha_{i'}+\beta_{j'}.
\]
Let $X_{\max}$ and $Y_{\max}$ be the respective maximizer sets.  The last
equality shows that $\mu$ maps $X_{\max}\times Y_{\max}$ into itself.
Minimality of $X'\times Y'$ implies
\[
    X_{\max}=X',
    \qquad
    Y_{\max}=Y'.
\]
Thus every $\alpha_i$ is equal and every $\beta_j$ is equal.  The kernel has
dimension two, proving the claim.

Choose $x+y-2$ independent rows, and delete one column from each side.  The
resulting square matrix has a nonzero determinant $D$: the row space is
precisely the direct sum of the two zero-sum coordinate spaces, for which
deleting one coordinate on each side is an isomorphism.  Modulo $p$, the two
side-constant vectors and the label vector
\[
    (a_i)_{i\in X'\cup Y'}
\]
all lie in the kernel.  The last vector is not a linear combination of the
first two.  Indeed, at least one side has two elements, since a one-pair
rectangle cannot be invariant under the fixed-point-free map $\mu$, and
distinct indices label distinct elements of $\mathbb F_p$.  Hence $p\mid D$.

On the $X'$-columns, each row is a difference $e_i-e_{i'}$ of two standard
basis vectors, and similarly on the $Y'$-columns.  Laplace expansion across
the $x-1$ retained $X'$-columns expresses $D$ as at most
\[
    \binom{x+y-2}{x-1}
\]
products of minors of matrices whose rows are differences of two standard
basis vectors.  Such matrices are totally unimodular, so every product is
zero or has absolute value one.  Therefore
\[
    p\le|D|\le\binom{x+y-2}{x-1}\le2^{x+y-2}.
\]
This is equivalent to $x+y\ge\log_2p+2$, and
$x\le|X|$, $y\le|Y|$ finish the proof.
\end{proof}

\section{A common random cut}

The following elementary hitting lemma is the probabilistic input in the
contraction.  We first introduce the objects that occur in its statement.

Let $V$ be a finite set.  A \emph{configuration on $V$} is a pair
$(j,\tau)$, where $j\in V$ and $\tau$ is a two-element multiset drawn from
$V$, with $j\notin\operatorname{supp}(\tau)$.  If $\tau=\{v,w\}$, define
\[
    \operatorname{supp}(\tau):=\{v,w\},
\]
regarded as an ordinary set; thus,
$\operatorname{supp}(\{v,v\})=\{v\}$.

A set $T\subseteq V$ \emph{hits} the configuration $(j,\tau)$ if
\[
    j\in T
    \qquad\text{and}\qquad
    T\cap\operatorname{supp}(\tau)\ne\varnothing.
\]
It hits a family of configurations if it hits at least one member of that
family.

\begin{lemma}[Common-cut lemma]
For every $\alpha>0$, there exist constants $D,m_0>0$, depending only on
$\alpha$, with the following property.  Let $V$ be a set of size
$m\ge m_0$, and let
\[
    \mathcal F=\{(j_t,\tau_t):1\le t\le s\}
\]
be a family of configurations on $V$, where $s\ge\alpha m$.  Suppose that
$j_1,\ldots,j_s$ are pairwise distinct and that
$\tau_1,\ldots,\tau_s$ are pairwise distinct.  Then there is a set
$T\subseteq V$, with
\[
    |T|\le D\sqrt m,
\]
that hits $\mathcal F$.

Moreover, suppose that
\[
    \mathcal F_1,\ldots,\mathcal F_q,
    \qquad q\le m,
\]
are families of configurations on the same set $V$, each satisfying the
preceding assumptions.  Then there is a single set $T\subseteq V$, with
$|T|\le D\sqrt m$, that hits all but at most $m/8$ of these families.
\end{lemma}

\begin{proof}
We first analyze one family
$\mathcal F=\{(j_t,\tau_t):1\le t\le s\}$.  Randomly partition $V$ as
\[
    V=P\sqcup Q.
\]
Retain a configuration $(j_t,\tau_t)$ if
\[
    j_t\in P
    \qquad\text{and}\qquad
    \operatorname{supp}(\tau_t)\subseteq Q.
\]
Because $j_t\notin\operatorname{supp}(\tau_t)$ and the support of $\tau_t$
has at most two elements, each configuration is retained with probability
at least $1/8$.  Hence the expected number of retained configurations is at
least
\[
    \frac{s}{8}\ge\frac{\alpha m}{8}.
\]
We may therefore fix a partition $V=P\sqcup Q$ for which the set
$\mathcal R$ of retained configurations satisfies
\[
    r:=|\mathcal R|\ge\frac{\alpha m}{8}.
\]

Let
\[
    \theta=\frac{C}{\sqrt m},
\]
where $C=C(\alpha)$ will be chosen later.  By increasing $m_0$, we may
assume that $\theta\le 1/2$.  Construct $T\subseteq V$ by including each
element independently with probability $\theta$.

First expose the choices defining $T\cap P$.  A retained configuration
$(j_t,\tau_t)$ is called \emph{activated} if $j_t\in T$.  Since the $j_t$
are pairwise distinct, the activation events are independent, each with
probability $\theta$.

Let $Z$ be the number of elements $v\in Q$ that belong to the support of at
least one activated configuration.  For each $v\in Q$, define
\[
    d_v:=
    \bigl|\{t:(j_t,\tau_t)\in\mathcal R,
                    \ v\in\operatorname{supp}(\tau_t)\}\bigr|.
\]
Thus, $d_v$ is the number of retained configurations whose multiset has
$v$ in its support.  Each of those configurations is activated
independently with probability $\theta$, and therefore
\[
    \mathbb P(v\text{ belongs to an activated support})
       =1-(1-\theta)^{d_v}.
\]
By linearity of expectation,
\[
    \mathbb E[Z]
      =\sum_{v\in Q}\bigl(1-(1-\theta)^{d_v}\bigr)
      \gg \sum_{v\in Q}\min\{1,\theta d_v\}.
\]

Set
\[
    H:=\{v\in Q:d_v\ge 1/\theta\}.
\]
By the definition of $H$, every $v\in H$ contributes $1$ to the last sum,
so
\[
    \mathbb E[Z]\gg |H|.
\]
If
\[
    |H|\ge\sqrt{\frac{\alpha m}{16}},
\]
then it follows that
\[
    \mathbb E[Z]\gg_\alpha\sqrt m.
\]

Suppose instead that
\[
    |H|<\sqrt{\frac{\alpha m}{16}}.
\]
The number of distinct two-element multisets supported entirely on $H$ is
\[
    \binom{|H|+1}{2}<\frac r2
\]
when $m_0$ is sufficiently large.  Since the multisets $\tau_t$ are
pairwise distinct, at least $r/2$ retained configurations have support
meeting $Q\setminus H$.  Consequently,
\[
    \sum_{v\in Q\setminus H}d_v\ge\frac r2.
\]
For $v\notin H$, we have $\theta d_v<1$, so
\[
    \mathbb E[Z]
      \gg \theta\sum_{v\in Q\setminus H}d_v
      \ge \frac{\theta r}{2}
      \gg \alpha C\sqrt m.
\]
Thus, in either case,
\[
    \mathbb E[Z]\gg_\alpha\sqrt m.
\]

We next bound the variance of $Z$.  For a retained configuration indexed by
$t$, let $Z^{(t,0)}$ denote the value of $Z$ obtained by forcing that
configuration not to be activated, while leaving all other activation
choices unchanged.  Put
\[
    \Delta_t:=Z-Z^{(t,0)}.
\]
Since $\operatorname{supp}(\tau_t)$ has at most two elements,
\[
    0\le\Delta_t\le 2.
\]
Moreover,
\[
    \sum_t\Delta_t\le Z,
\]
because an element counted by $Z$ contributes to $\Delta_t$ only when $t$
is the unique activated configuration whose support contains that element.
The conditional-variance form of the Efron--Stein inequality therefore gives
\[
    \operatorname{Var}(Z)
      \le \mathbb E\!\left[\sum_t\Delta_t^2\right]
      \le 2\mathbb E\!\left[\sum_t\Delta_t\right]
      \le 2\mathbb E[Z].
\]
Hence, by Chebyshev's inequality,
\[
    \mathbb P\!\left(Z<\frac12\mathbb E[Z]\right)
      \le \frac{8}{\mathbb E[Z]}
      =O_\alpha(m^{-1/2}).
\]

Conditional on the choices in $P$, the choices defining $T\cap Q$ remain
independent.  If
\[
    Z\ge\frac12\mathbb E[Z]\gg_\alpha\sqrt m,
\]
the probability that $T\cap Q$ contains none of these $Z$ elements is at
most
\[
    (1-\theta)^Z
      \le \exp(-\theta Z)
      \le \exp(-c_\alpha C).
\]
If $T\cap Q$ contains one of them, then some activated configuration
$(j_t,\tau_t)$ satisfies
\[
    j_t\in T
    \qquad\text{and}\qquad
    T\cap\operatorname{supp}(\tau_t)\ne\varnothing.
\]
Thus $T$ hits $\mathcal F$.  The probability that $T$ fails to hit
$\mathcal F$ is therefore at most
\[
    \delta(C,m):=\exp(-c_\alpha C)+O_\alpha(m^{-1/2}).
\]
Choose $C$, and then $m_0$, sufficiently large that
$\delta(C,m)\le 1/64$.

Also,
\[
    \mathbb E[|T|]=\theta m=C\sqrt m.
\]
By Markov's inequality,
\[
    \mathbb P(|T|>4C\sqrt m)\le\frac14.
\]
Together with the preceding failure bound, this proves the assertion for
one family.

Finally, suppose that at most $m$ families are specified on the same set
$V$.  For each family, fix its own partition $V=P\sqcup Q$ retaining at
least one eighth of its configurations.  Now construct one common random
set $T$ by including each element of $V$ independently with probability
$\theta$.

The preceding analysis applies to each family and shows that its probability
of being missed by $T$ is at most $1/64$.  No independence between the
different failure events is required.  If $M$ denotes the number of missed
families, then linearity of expectation gives
\[
    \mathbb E[M]\le\frac m{64}.
\]
Markov's inequality gives
\[
    \mathbb P\!\left(M>\frac m8\right)\le\frac18.
\]
At the same time,
\[
    \mathbb P(|T|>4C\sqrt m)\le\frac14.
\]
Therefore, with positive probability, both
\[
    M\le\frac m8
    \qquad\text{and}\qquad
    |T|\le 4C\sqrt m
\]
hold.  Taking $D=4C$ completes the proof.
\end{proof}

The next lemma supplies the combinatorial contraction.  The root, when
present, is a single additional vertex shared by all rooted pieces.  For
this reason the correct metric bound is a radius bound from the root, not a
diameter-four bound on the whole rooted component.

\begin{lemma}[Bounded-radius contraction]
Let $G$ be a graph with $N$ designated live vertices and, optionally, one
additional vertex $0$, called the root.  An edge between two live vertices
is called a live edge, and an edge joining a live vertex to $0$ is called a
root edge.  Suppose that at most $f$ live vertices are incident with neither
a live edge nor a root edge.

Then $G$ contains a forest $F$ with the following properties:
\begin{enumerate}
    \item every component of $F$ not containing $0$ has diameter at most
          three;
    \item every live vertex in the component containing $0$ has distance at
          most four from $0$, so that this component has diameter at most
          eight;
    \item if one live vertex is retained from each component not containing
          $0$, and no live vertex is retained from the component containing
          $0$, then at most
          \[
              \frac{N+f}{2}
          \]
          live vertices are retained.
\end{enumerate}
\end{lemma}

\begin{proof}
Let $L$ be the set of live vertices, and let
\[
    G_{\mathrm{live}}:=G[L]
\]
be the subgraph induced by $L$.  Thus $G_{\mathrm{live}}$ contains all live
vertices and all edges joining two live vertices, but it contains neither
the root nor any root edge.

Take a maximal matching $M$ in $G_{\mathrm{live}}$.  We construct a spanning
forest $F\subseteq G$, beginning with all the edges of $M$.  Every unmatched
vertex $u$ that has a neighbor in $G_{\mathrm{live}}$ has a matched neighbor.
Indeed, if $u$ had an unmatched neighbor $v$, then the edge $\{u,v\}$ could
be added to $M$, contradicting maximality.

For each such $u$, choose a matched neighbor $v$.  There is a unique matching
edge $e\in M$ containing $v$.  Assign $u$ to $e$, and add the edge $\{u,v\}$
to $F$.  For each matching edge $e=\{a,b\}$, the edge $e$, together with the
vertices assigned to it and their chosen edges, forms a double-star.  These
double-stars are pairwise disjoint and have diameter at most three.

If a double-star contains a live vertex joined to the root $0$, add exactly
one such root edge to $F$.  Because only one root edge is added from each
double-star, this creates no cycle.  All double-stars attached to $0$ become
part of the same rooted component.  If $x$ is the attachment point of one
such double-star, every vertex in that double-star is at distance at most
three from $x$.  Therefore every live vertex in the rooted component is at
distance at most four from $0$.  Thus the rooted component has radius at
most four from $0$, and hence diameter at most eight.

Now consider an unmatched vertex that is isolated in
$G_{\mathrm{live}}$.  If it has a root edge, add one such edge to $F$.
Otherwise, leave it as a singleton component.  Let $f_0$ be the number of
these exceptional singleton vertices.  Each is incident with neither a live
edge nor a root edge, so the hypothesis gives
\[
    f_0\le f.
\]

Every component of $F$ not containing $0$ is either an unrooted double-star,
of diameter at most three, or an exceptional singleton, of diameter zero.
Each unrooted double-star contains at least two live vertices and contributes
one retained vertex.  The component containing $0$ contributes none, while
the $f_0$ exceptional singleton components contribute $f_0$.  Hence the
total number of retained live vertices is at most
\[
    \frac{N-f_0}{2}+f_0
      =\frac{N+f_0}{2}
      \le\frac{N+f}{2}.
\]
This proves all three assertions.
\end{proof}

\section{Integral contraction}

We now construct many primitive lattice pivots.  All norms below are
$\ell_1$-norms.  Put
\[
    L:=\left\lfloor\frac{\log_2 p}{3}\right\rfloor.
\]
We may assume that $p$ is large enough that $L\ge 2$.

\begin{proposition}[Compression]
There are an integer $K$ and a subgroup $P\subseteq\Lambda$, primitive in
$\mathbb Z^n$, such that
\[
    \operatorname{rank}P=n-K,
    \qquad
    K=O\!\left(\sqrt n+\frac{n}{\log p}\right),
\]
and $\mathbb Z^n/P$ is free of rank $K$.  Moreover, the quotient has a basis
$c_1,\ldots,c_K$ in which the image of every original coordinate $e_a$ has
$\ell_1$-norm at most $n^{O(1)}$.
\end{proposition}

\begin{proof}
We construct an increasing sequence of primitive subgroups
\[
    P_0\subseteq P_1\subseteq\cdots\subseteq\Lambda.
\]
Writing $Q_r:=\mathbb Z^n/P_r$, enlarging $P_r$ imposes additional relations
and therefore eliminates basis directions from the quotient:
\[
    Q_{r+1}\cong Q_r/(P_{r+1}/P_r).
\]
Alongside this algebraic operation, we keep track of a convenient basis of
each $Q_r$, dividing its elements into core coordinates and live
representatives.  Promoting a live representative to the core does not
change the quotient; it only declares that basis element permanent.  The
quotient changes only when new relations are added to $P_r$.

The construction keeps some quotient-basis elements permanently; we call
these the \emph{core coordinates}.  The remaining original indices are
grouped into pairwise disjoint \emph{live blocks}.  A live block is a set
$B\subseteq[n]$ equipped with a quotient-basis element $x_B$ such that every
$a\in B$ has image
\[
    \overline e_a=x_B+z_a,
\]
where $z_a$ is supported on the core coordinates.  The representative $x_B$
is still eligible to be eliminated through a later forest contraction, and
the weight of the block is $|B|$.  When $x_B$ is permanently retained as a
core coordinate, we call $B$ nonlive.

At the beginning of round $r$, let $k_r$ be the number of core coordinates
and let $m_r$ be the number of live blocks.  We maintain the following
invariant: the quotient $Q_r$ is free on a basis consisting of the $k_r$
core coordinates and one live
representative $x_B$ for each of $m_r$ pairwise disjoint live blocks
$B\subseteq[n]$.

For every original index $a\in[n]$, its image $\overline e_a$ in $Q_r$ has
the form
\[
    \overline e_a=
    \begin{cases}
        x_B+z_a, & a\in B\text{ for a live block }B,\\
        z_a,     & a\text{ belongs to a nonlive block},
    \end{cases}
    \qquad
    \|z_a\|_1\le M_r,
\]
where every $z_a$ is supported on the core coordinates.  Initially
\[
    P_0=\{0\},\qquad k_0=0,\qquad m_0=n,\qquad M_0=1,
\]
and the live blocks are the singletons.  Every live block has weight
$|B|<L$.  Whenever a block reaches weight at least $L$, we promote its
representative to the core and declare the block nonlive.

Fix one round and suppress the subscript $r$, writing $k,m,M$ for the current
parameters.  Label the live blocks $B_1,\ldots,B_m$ and write $x_i$ for the
representative of $B_i$.

For distinct live blocks $B_i,B_j$, we have
\[
    |B_i|+|B_j|<2L<\log_2p+2.
\]
By Lemma~2.1, the rectangle formed by the index pairs with one index in
$B_i$ and the other in $B_j$ cannot be $\mu$-invariant.  For every ordered
pair $i\ne j$, choose a departing index pair $t_{ij}$ in this rectangle such
that $\mu(t_{ij})$ lies outside it.  Project all the corresponding collision
rows into the quotient basis fixed at the beginning of the round.

For a multiset $\nu$ of live labels, write
\[
    x(\nu):=\sum_{a\in\nu}x_a,
\]
with multiplicity.  The projected departing row has the form
\begin{equation}
    r_{ij}=x_i+x_j-x(\nu_{ij})+g_{ij},
    \qquad
    |\nu_{ij}|\le2,
    \qquad
    \nu_{ij}\ne\{i,j\},
    \qquad
    \|g_{ij}\|_1\le4M,
    \tag{4.1}
    \label{eq:departing-row}
\end{equation}
where $g_{ij}$ is supported on the core.  Indeed, a collision row is the
signed sum of four original coordinate images.  The inequality
$\nu_{ij}\ne\{i,j\}$ follows from the fact that $t_{ij}$ departs from the
rectangle.

Define the \emph{unit graph} for the round to have vertex set
$\{0,1,\ldots,m\}$, where the vertices $1,\ldots,m$ label the live
representatives and $0$ is an auxiliary root representing the core.  Call a
projected row a \emph{unit row} if, after multiplying the row by $-1$ if
necessary, its noncore part is either $x_a$ or $x_a-x_b$.  In the first case
attach the row to the root edge $\{0,a\}$; in the second case attach it to
the edge $\{a,b\}$.  The graph records only the noncore part of each row;
the core-supported term is retained as part of the row attached to the
edge.  If several rows produce the same edge, retain any one of them.  Begin
the unit graph with the edges coming from those projected departing rows
$r_{ij}$ that are already unit rows.

Fix an anchor $i$ and vary the label $j\ne i$.  Classify the row in
\eqref{eq:departing-row} as follows.
\begin{itemize}
    \item If $\nu_{ij}$ contains $j$, cancellation leaves a unit row incident
          with $i$.  Call $i$ \emph{direct} if this happens for at least one
          label $j$.
    \item If $\nu_{ij}$ contains $i$, cancellation leaves a unit row incident
          with $j$, so $j$ is covered by the unit graph.
    \item If $\nu_{ij}$ contains neither $i$ nor $j$ and
          $|\nu_{ij}|\le1$, call $j$ \emph{easy for $i$}.
    \item In the remaining case, $\nu_{ij}$ is a two-element multiset
          disjoint from $\{i,j\}$.  Call $j$ \emph{proper for $i$}, with
          type $\nu_{ij}$.
\end{itemize}

For a fixed anchor $i$, group its proper labels by type.  Within each group,
choose one label $j_0$ and subtract its row from the other rows in that
group.  If $j$ has the same type as $j_0$, this gives
\[
    r_{ij}-r_{ij_0}
      =x_j-x_{j_0}+(g_{ij}-g_{ij_0}),
\]
which is a unit row whose core part has norm at most $8M$.  Add its edge to
the unit graph.  For each proper type, the resulting edges form a star
centered at $j_0$ and cover all but at most one proper label of that type.

Let $a_i$ be the number of easy labels for anchor $i$, and let $s_i$ be the
number of its distinct proper types.  Suppose that some nondirect anchor $i$
satisfies
\[
    a_i<\frac m8,
    \qquad
    s_i<\frac m8.
\]
The immediate unit rows and the proper-type star rows then cover at least
\[
    m-1-a_i-s_i>\frac{3m}{4}-1
\]
live labels.  Thus the unit graph has at most $m/4+1$ isolated live
vertices.  Lemma~3.2 contracts this graph to at most
\[
    \frac{m+(m/4+1)}{2}=\frac{5m}{8}+\frac12
\]
live blocks.

It remains to treat the alternative in which no nondirect anchor has both
$a_i<m/8$ and $s_i<m/8$.  Partition the nondirect anchors into two classes:
declare $i$ to be an easy anchor if $a_i\ge m/8$, and otherwise declare it
to be a proper anchor, in which case necessarily $s_i\ge m/8$.

For each easy anchor $i$, let $E_i$ be its set of easy labels, so
$|E_i|\ge m/8$.  Include each live label independently in a random set with
probability $\theta=C/\sqrt m$.  For a fixed easy anchor,
\[
    \mathbb P(T\cap E_i=\varnothing)
      =(1-\theta)^{|E_i|}
      \le \exp(-C\sqrt m/8).
\]
A union bound over at most $m$ easy anchors shows that, for all sufficiently
large $m$, the probability of missing any $E_i$ is less than $1/4$.  Since
the expected sample size is $C\sqrt m$, Markov's inequality gives
\[
    \mathbb P(|T|>2C\sqrt m)\le\frac12.
\]
Hence there is a deterministic set $T_{\mathrm{easy}}$ of size at most
$2C\sqrt m$ that meets every easy-label set $E_i$.

For each proper anchor $i$, choose one proper label from each of at least
$m/8$ distinct proper types.  The resulting pairs
\[
    (j,\nu_{ij})
\]
form a family of configurations on the set of live labels: the chosen labels
are distinct, the types are distinct, and $j$ does not belong to the support
of $\nu_{ij}$.  Applying Lemma~3.1 simultaneously to these at most $m$
families, with $\alpha=1/8$, gives a set $T_{\mathrm{prop}}$ of size
$O(\sqrt m)$ that hits all but at most $m/8$ proper-anchor families.

Set
\[
    T:=T_{\mathrm{easy}}\cup T_{\mathrm{prop}}.
\]
Then $|T|=O(\sqrt m)$, and this single set has both required properties.
Promote all representatives indexed by $T$ to core coordinates.  If an easy
anchor $i$ does not itself lie in $T$, choose an easy label $j\in E_i\cap T$;
then \eqref{eq:departing-row} becomes a unit row incident with $i$.  If a proper anchor
$i\notin T$ belongs to a family hit by $T_{\mathrm{prop}}$, then some chosen
label $j$ and at least one element of the support of $\nu_{ij}$ lie in $T$;
again \eqref{eq:departing-row} becomes a unit row incident with $i$.

Direct rows and all pre-existing unit rows survive the promotion: if one
endpoint of a unit edge enters the core, that edge simply becomes a root
edge.  Consequently, among the remaining live labels, at most $m/8$ are
isolated.  Lemma~3.2 therefore leaves at most
\[
    \frac{m+m/8}{2}=\frac{9m}{16}
\]
live blocks.

In either case, once $m$ is above a fixed threshold $m_*$, the number of live
blocks falls by a factor at most $2/3$.  Choose the forest rows supplied by
Lemma~3.2.  Eliminate every live representative in the rooted component and
all but one representative in each unrooted component.  Order the selected
rows from the leaves toward the retained representative or toward the root.
Each row then has coefficient $1$ on a fresh pivot coordinate.  The pivot
matrix is unit triangular, so the selected rows extend $P_r$ primitively.
The retained live block is the union of the physical blocks in its forest
component.

We next record coefficient growth.  Before promotion, a raw unit row has
core norm at most $4M$, and a proper-type difference has core norm at most
$8M$.  Reclassifying sampled representatives as core coordinates introduces
at most four further unit coefficients, so every selected unit row has core
norm at most $12M$.  Promotion changes the core-norm bound on each original
coordinate from $M$ to at most $M+1\le2M$.

The corrected metric conclusion of Lemma~3.2 gives every eliminated live
representative a substitution path of length at most four, either to the
retained representative of its unrooted component or to the root.  Substituting
along such a path changes the bound on an original coordinate to at most
\[
    2M+4\cdot12M=50M.
\]
Promoting a newly heavy representative costs at most one further unit
coefficient.  Thus the next round may take
\[
    M_{r+1}\le60M_r.
\]
There are only $O(\log n)$ nonterminal rounds, and hence throughout the
construction
\[
    M_r\le60^{O(\log n)}=n^{O(1)}.
\]

It remains to count the final core coordinates.  The common samples
contribute
\[
    O\!\left(\sum_{\ell\ge0}\sqrt{n(2/3)^\ell}\right)
      =O(\sqrt n)
\]
coordinates.  Every block promoted for being heavy has weight at least $L$,
and these physical blocks are pairwise disjoint, so heavy promotions
contribute at most
\[
    \frac nL=O\!\left(\frac n{\log p}\right)
\]
coordinates.  Once fewer than $m_*$ live blocks remain, promote all their
representatives; this contributes $O(1)$ further coordinates.

Let $R$ be the terminal stage and define
\[
    P:=P_R,
    \qquad
    K:=k_R.
\]
The final quotient is free on its $K$ core coordinates, and
\[
    K=O\!\left(\sqrt n+\frac n{\log p}\right).
\]
Since its rank is $K$, exactly $n-K$ primitive pivots were selected, so
$\operatorname{rank}P=n-K$.  The coefficient-growth estimate shows that the
image of every original coordinate has $\ell_1$-norm at most $n^{O(1)}$ in
the final core basis.  This completes the proof.
\end{proof}

\section{The final determinant}

\begin{proof}[Proof of Theorem 1.2]
Let $P\subseteq\Lambda$ be the rank-$(n-K)$ primitive subgroup from
Proposition~4.1.  Since $\Lambda$ has rank $n-1$ by Lemma~1.3, choose
$K-1$ original collision rows whose images modulo $P$ are rationally
independent.  In a unimodular coordinate system adapted to the primitive
pivots, eliminate their pivot-coordinate entries using the selected rows.
The remaining $(K-1)$-by-$K$ matrix consists of their images in the core.

Every original coordinate has core norm at most $M=n^{O(1)}$, and an
original collision row is the signed sum of four such images.  Hence every
row of this core matrix has norm at most $4M$.  Select a nonzero
$(K-1)$-minor $D$.  It is a maximal minor of the full $(n-1)$-row lattice
matrix after unimodular row and column operations.  Lemma~1.3 therefore
gives $p\mid D$.  Hadamard's inequality gives
\[
    p\le|D|\le(4M)^{K-1}=\exp\bigl(O(K\log n)\bigr).
\]
Writing $X=\log p$ and using the bound for $K$, we obtain
\[
    X\le C\left(\sqrt n+\frac nX\right)\log n
\]
for an absolute constant $C$.  If $n\ge X^2$, the asserted lower bound is
immediate.  Otherwise
\[
    \frac nX\le\sqrt n
    \qquad\text{and}\qquad
    \log n\le2\log X,
\]
whence
\[
    X\ll\sqrt n\log X.
\]
Rearranging yields
\[
    n\gg
    \left(\frac{\log p}{\log\log p}\right)^2,
\]
as required.
\end{proof}

\begin{corollary}
As $p\to\infty$,
\[
    m(p)=(\log p)^{2+o(1)}.
\]
\end{corollary}

\begin{proof}
Theorem~1.2 gives
\[
    m(p)\ge(\log p)^2(\log\log p)^{-2}
          =(\log p)^{2-o(1)}.
\]
Bedert's construction gives $m(p)\ll(\log p)^2$~\cite{bedert}.
\end{proof}

\begin{thebibliography}{9}

\bibitem{bedert}
B.~Bedert,
\emph{On unique sums in Abelian groups},
Combinatorica \textbf{44} (2024), 269--298;
\url{https://doi.org/10.1007/s00493-023-00069-w}.

\bibitem{green}
B.~Green,
\emph{100 open problems}, Problem~27;
\url{https://people.maths.ox.ac.uk/greenbj/papers/open-problems.pdf}.

\end{thebibliography}

\end{document}

